\ifdefined\iscombined
	\lecturetitle{Introduction}
\else
\documentclass{article}
\usepackage{xparse}
\usepackage{changepage}
\usepackage{listings}
\usepackage{amsthm}
\usepackage{comment}
\usepackage{amsmath}
\usepackage{braket}
\usepackage{amsfonts}
\usepackage{sectsty}
\usepackage{hyperref}
\usepackage{fancyhdr}
\usepackage[top=1in,bottom=1in,left=1in,right=1in]{geometry}
\usepackage{float}
\usepackage{graphicx}
\usepackage{mathtools}

\date{
	\vspace{-.75em}
	{\large \today}
}

\title{
	\vspace*{-1.5em}
	{\Huge Lecture 1} \\
	\vspace{-1em}
	\rule{\textwidth/4}{.01em} \\
	\vspace{-.25em}
	{\Large Introduction}
	\vspace{-.75em}
}

\author{
	{\large Matthew Farah}
}

\hypersetup{
	colorlinks=true,
	linkcolor=blue,
	filecolor=magenta,
	urlcolor=blue,
	citecolor=blue,
	anchorcolor=blue
}

\fancyhead{}
\fancyhead[L]{\today}
\fancyhead[R]{Lecture 1 - Introduction}

\setlength{\parindent}{0cm}

\newcounter{example}
\renewcommand{\theexample}{\arabic{example}}

\newenvironment{example}[1][]{
	\refstepcounter{example}
	\begin{adjustwidth}{1.5em}{}
	\noindent\textbf{\ifx&#1&Example \theexample:\else Example \theexample \ (#1):\fi}
	\ignorespaces
}{
	\vspace{.5em}
	\end{adjustwidth}
}

\newenvironment{solution}{
	\vspace{.5em}
	\par
	\begin{adjustwidth}{1.5em}{}
	\textbf{{Solution:}}
	\par
	\vspace{.5em}
}{
	\vspace{1em}
	\end{adjustwidth}
}

\newcounter{subexample}
\renewcommand{\thesubexample}{\alph{subexample}}

\newenvironment{subexample}{
	\setcounter{step}{0}
	\refstepcounter{subexample}
	\vspace{1em}\par\noindent
	\begin{adjustwidth}{1.5em}{}
	\noindent\textbf{(\thesubexample)}
	\ignorespaces
}{
	\par
	\vspace{1em}
	\end{adjustwidth}
}

\renewenvironment{proof}{
	\vspace{.5em}\par\noindent
	\begin{adjustwidth}{1.5em}{}
	\emph{Proof:}\par\vspace{1em}
	\setlength{\parindent}{0cm}
}{
	\end{adjustwidth}
}

\newtheorem{theorem}{Theorem}

\renewenvironment{theorem}[1][]{
	\refstepcounter{theorem}
	\vspace{1em}\par\noindent
	\begin{adjustwidth}{1.5em}{}
	\noindent\textbf{\ifx&#1&Theorem \thetheorem:\else Theorem \thetheorem \ (#1):\fi}
	\ignorespaces
}{
	\vspace{.5em}
	\end{adjustwidth}
}

\newtheorem{lemma}{Lemma}

\renewenvironment{lemma}[1][]{
	\refstepcounter{lemma}
	\vspace{1em}\par\noindent
	\begin{adjustwidth}{1.5em}{}
	\noindent\textbf{\ifx&#1&Lemma \thelemma:\else Lemma \thelemma \ (#1):\fi}
	\ignorespaces
}{
	\vspace{.5em}
	\end{adjustwidth}
}

\newcounter{definition}
\renewcommand{\thedefinition}{\arabic{definition}}

\newenvironment{definition}[1][]{
	\refstepcounter{definition}
	\vspace{1em}\par\noindent
	\begin{adjustwidth}{1.5em}{}
	\noindent\textbf{\ifx&#1&Definition \thedefinition:\else Definition \thedefinition \ (#1):\fi}
	\ignorespaces
}{
	\vspace{.5em}
	\end{adjustwidth}
}

\newcounter{corollary}
\renewcommand{\thecorollary}{\arabic{corollary}}

\newenvironment{corollary}[1][]{
	\refstepcounter{corollary}
	\vspace{1em}\par\noindent
	\begin{adjustwidth}{1.5em}{}
	\noindent\textbf{\ifx&#1&Corollary \thecorollary:\else Corollary \thecorollary \ (#1):\fi}
	\ignorespaces
}{
	\vspace{.5em}
	\end{adjustwidth}
}

\newcounter{proposition}
\renewcommand{\theproposition}{\arabic{proposition}}

\newenvironment{proposition}[1][]{
	\refstepcounter{proposition}
	\vspace{1em}\par\noindent
	\begin{adjustwidth}{1.5em}{}
	\noindent\textbf{\ifx&#1&Proposition \theproposition:\else Proposition \theproposition \ (#1):\fi}
	\ignorespaces
}{
	\vspace{.5em}
	\end{adjustwidth}
}

\newenvironment{remark}{
	\vspace{.5em}\par\noindent
	\begin{adjustwidth}{1.5em}{}
	\noindent\textbf{Remark:}
	\ignorespaces
}{
	\vspace{.5em}
	\end{adjustwidth}
}

\makeatother
\makeatletter

\newcounter{step}
\renewcommand{\thestep}{\Roman{step}}

\newenvironment{step}{
	\refstepcounter{step}
	\vspace{1em}\par\noindent
	\ignorespaces\noindent\textbf{(\thestep)}
	\ignorespaces
}{
	\par
}

\sectionfont{\fontsize{12}{12}\bfseries}
\subsectionfont{\fontsize{10}{12}\normalfont}
\subsubsectionfont{\fontsize{10}{12}\normalfont}

\begin{document}

\maketitle
\pagestyle{fancy}

\fi

\begin{definition}
	A set $\mathbb{F}$ is called a field if $\mathbb{F}$ has:
	\begin{enumerate}
		\item A zero element (denoted 0).
		\item A one element (denoted 1).
		\item A defined notion of addition on elements of the set.
		\item A defined notion of multiplication on elements of the set. \\
	\end{enumerate}
	
	The addition and multiplication on $\mathbb{F}$ must satisfy, for all $a, b, c \in \mathbb{F}$:
	\begin{enumerate}
		\item Commutativity.
		\begin{itemize}
			\item $a + b = b + a$
			\item $ab = ba$.
		\end{itemize}
		\item Associativity.
		\begin{itemize}
			\item $(a + b) + c = a + (b + c)$ 
			\item $(ab)c = a(bc)$.
		\end{itemize}
		\item Distributivity over addition.
		\begin{itemize}
			\item $a(b + c) = ab + ac$. \\
		\end{itemize}
	\end{enumerate}

	As well, addition and multiplication on $\mathbb{F}$ must:
	\begin{enumerate}
		\item Have an identity element (0 and 1 respectively) such that for all $a \in \mathbb{F}$.
		\begin{itemize}
			\item $a + 0 = 0 + a = a$
			\item $1 \cdot a = a \cdot 1 = a$.
		\end{itemize}
		\item Have an inverse element.
		\begin{itemize}
			\item $\forall a \in \mathbb{F}, \exists b \in \mathbb{F}$ such that $a + b = 0$
			\item $\forall a \in \mathbb{F} \setminus \set{0}, \exists b \in \mathbb{F}$ such that $a b = 1$
		\end{itemize}
	\end{enumerate}
\end{definition}

\begin{example}
	Is the set of all natural numbers ($\mathbb{N}$) a field?
\end{example}

\begin{solution}
	No! There does not exist an element $a \in \mathbb{N}$ such that $1 + a = 0$. In other words, no elements besides 0 in $\mathbb{N}$ have an additive inverse.
\end{solution}

\begin{example}
	Is the set of all integers ($\mathbb{Z}$) a field?
\end{example}

\begin{solution}
	No! There does not exist an element $a \in \mathbb{Z}$ such that $2a = 1$. In other words, no elements besides 1 in $\mathbb{Z}$ have a multiplicative inverse.
\end{solution}

\begin{example}
	Is the set of all rational numbers ($\mathbb{Q}$) a field?
\end{example}

%TODO: Give proper proof of this.
\begin{solution}
	Yes!
\end{solution}

\begin{example}
	Is the set of all real numbers ($\mathbb{R}$) a field?
\end{example}

\begin{solution}
	Yes! Although the proof is quite complicated. In many undergraduate courses, $\mathbb{R}$ is often defined axiomatically as a complete, ordered field.
\end{solution}

\begin{example}
	Is the set of all complex numbers ($\mathbb{C}$) a field?
\end{example}

%TODO: Give proper proof of this.
\begin{solution}
	Yes!
\end{solution}

\ifdefined\iscombined
	\newpage
\else
	\end{document}
\fi
